\documentclass[11pt]{article}
\usepackage[a4paper,margin=1in]{geometry}
\usepackage{amsmath,amssymb,booktabs,array,microtype}
\usepackage[colorlinks=true,linkcolor=blue,urlcolor=blue,citecolor=blue]{hyperref}
\newcommand{\R}{\mathbb R}
\newcommand{\Q}{\mathbb Q}
\newcommand{\sig}{\sigma}
\title{The Salem boundary for substitution spectral measures:\\
scoped later-literature answers to arXiv:1305.7373}
\author{Literature status note --- agent\_lane\_12}
\date{22 July 2026}

\begin{document}
\maketitle

\noindent\textbf{Status.}
\emph{Literature-already-answered partial subcases; the exact log-H\"older
bound remains open.}

\section*{Original question}
Bufetov and Solomyak \cite{BS14} consider a primitive aperiodic substitution
with irreducible substitution matrix and Salem Perron--Frobenius eigenvalue
\(\alpha\).  On page 15 in Section 4, immediately after Theorem 4.1, they say that they
do not know whether almost-every suspension flow has H\"older spectral
estimates in this borderline case and expect failure.  In Section 5,
immediately after Theorem 5.1 and equation (5.1), they ask whether the
self-similar flow nevertheless satisfies (page 26, Remark 5.3)
\[
 \sig_a([\omega-r,\omega+r])
 \le C_B(\log(1/r))^{-\gamma},
 \qquad |\omega|\in[B^{-1},B],\quad r\downarrow0. \tag{Q}
\]

\section*{Explicit later answers}
Marshall-Maldonado \cite{MM24} states that the main objective is the Salem
question raised in \cite{BS14}.  Theorems 1.1--1.2 prove H\"older local bounds
for the self-similar flow at algebraic spectral parameters
\(\omega\in\Q(\alpha)\).  The exponent is controlled by arithmetic parameters,
but the admissible starting radius depends on \(\omega\); hence this is not the
uniform estimate (Q).

Marshall-Maldonado and Solomyak \cite{MS26} give the following current
boundary.  Their results concern Lip-cylindrical mean-zero functions; the
letter-cylinder measures in \cite{BS14} are included after subtracting the
constant component away from frequency zero.

\begin{center}
\begin{tabular}{>{\raggedright\arraybackslash}p{0.19\textwidth}
                >{\raggedright\arraybackslash}p{0.34\textwidth}
                >{\raggedright\arraybackslash}p{0.36\textwidth}}
\toprule
Result in \cite{MS26} & Scope & Conclusion \\
\midrule
Theorem 2.3 (pp. 6--7) & Lebesgue-a.e. roof vector & Uniform spectral modulus
\(\exp[-\beta L_3/L_4]\), where \(L_3\) and \(L_4\) are the relevant third
and fourth iterated logarithms; quantitative weak mixing follows. \\
Theorem 2.4 (p. 7) & Every roof vector; almost every Lip-cylindrical test function &
Every uniform H\"older bound \(Cr^\eta\), \(\eta>0\), fails along a sequence
of scales.  This confirms the expected failure of the Section 4 H\"older
regime. \\
Theorem A.2 (p. 18) & The self-similar roof & Uniformly in \(\omega\in\R\),
\(\sig_f(B_r(\omega))\le a_\alpha\exp[-c_\alpha\log^*(r^{-1})]\), with the
same quantitative weak-mixing rate. \\
\bottomrule
\end{tabular}
\end{center}

\section*{What remains open}
Theorem A.2 is the first uniform quantitative bound for the self-similar Salem
case, but \(\exp[-c\log^*(1/r)]\) is much weaker than every negative power of
\(\log(1/r)\).  The authors of \cite{MS26} explicitly say that a log-H\"older
rate seems plausible and is not known.  Thus the later literature answers
substantial parts of the extracted Salem question but neither proves nor
disproves (Q).

The arithmetic obstruction is visible in \cite[Proposition A.3]{MS26}: the
escape scales satisfy \(k_{i+1}=L^{k_i}\), giving only
\(\sum_{n<N}\|w\alpha^n\|_{\R/\mathbb Z}^2\gtrsim\log^*N\).  A proof of (Q)
would require a qualitatively denser uniform escape estimate.

\section*{Provenance and search evidence}
Both supporting papers explicitly cite \cite{BS14} as the source of the Salem
question, so this is a direct literature answer rather than an agent-inferred
reformulation.  Searches covered the exact source sentence, ``Salem
substitution suspension flow spectral measure'', the three arXiv identifiers,
and the cited 2025 discrete-action sequel.  No later proof or disproof of (Q)
was found.

\begin{thebibliography}{9}
\bibitem{BS14}
A.~I. Bufetov and B.~Solomyak,
\emph{On the modulus of continuity for spectral measures in substitution dynamics},
Adv. Math. 260 (2014), 84--129; arXiv:1305.7373.

\bibitem{MM24}
J.~Marshall-Maldonado,
\emph{Modulus of continuity for spectral measures of suspension flows over Salem type substitutions},
Israel J. Math. 263 (2024), 657--685; arXiv:2009.13607.

\bibitem{MS26}
J.~Marshall-Maldonado and B.~Solomyak,
\emph{Quantitative weak mixing for typical Salem substitution suspension flows},
arXiv:2601.15035 (2026), Theorems 2.3, 2.4, and A.2.
\end{thebibliography}

\end{document}
